% !TEX root = ../nc-quadric-classification.tex
% 2026-11-21 / v1.11
\section{Artin 基底での foundation と反 Serre 偏極}
\label{sec:artin-foundation}

この節の基底は、剰余体 $k$ を持つ有限次元可換局所 Artin
$k$-代数 $(A,\mathfrak m,k)$ に限る。
非被約な一般有限型基底への代数化は、この節の結論には含めない。
$(\mathcal C,\mathcal O)\in\HH$ とし、
$\mathcal G=\operatorname{Ind}\mathcal C$ と置く。
flat Abel 変形は Lowen--Van den Bergh の意味で取り、
Grothendieck 圏で記すときは $\mathcal G_A$ とする。

\begin{theorem}[Artin 変形への foundation の持上げ]
\label{thm:artin-foundation-lift}
定理~\ref{thm:normal-helix} の一方の helix $(E_i)_{i\in\mathbb Z}$ を選び、
$E_0=\mathcal O$、$T=\bigoplus_{i=0}^{3}E_i$、
$\Lambda=\operatorname{End}(T)$ と置く。
任意の flat Grothendieck 変形 $\mathcal G_A$ において、次が成立する。
\begin{enumerate}
\item 各 $E_i$ は flat な noetherian 対象 $E_{i,A}$ に持ち上がり、
  その持上げの同型類は一意である。同型そのものの一意性は主張しない。
\item 持上げた列は geometric helix であり、
  $T_A=\bigoplus_{i=0}^{3}E_{i,A}$ は $D(\mathcal G_A)$ の compact tilting generator である。
\item $\Lambda_A=\operatorname{End}(T_A)$ は階数二十四の自由 $A$-代数で、
  $\Lambda_A\otimes_A k\simeq\Lambda$。
  これは補題~\ref{lem:foundation-artin-deformation} の有限 foundation 表示に入る。
\item $D(\mathcal G_A)\simeq D(\operatorname{Mod}\Lambda_A)$、
  $D(\mathcal G_A)^c\simeq\operatorname{Perf}\Lambda_A$ である。
  compact 部分は $A$ 上 smooth proper であり、相対 Serre 関手を持つ。
\end{enumerate}
さらに、閉ファイバーを同定した Abel 変形と有限 foundation 代数の線形変形は、
De Deken--Lowen の変形比較により対応する。
この対応は適合平方根または AS4 被覆の持上げを主張しない。
\end{theorem}
\begin{proof}
完全な証明は相対モジュライ論文 \cite[\texttt{artin_foundation}]{QuadricModuli2026} に置く。
\end{proof}

\begin{lemma}[Artin 基底での exactness の認識]
\label{lem:artin-morita-heart-exact}
定理~\ref{thm:artin-foundation-lift} の表示の下で、$P$ を
$A$-線形の可逆 perfect $\Lambda_A$-双加群とし、
$F_A=-\otimes_{\Lambda_A}^{\mathbf L}P$ を $D(\mathcal G_A)$ の自己同値とする。
その特殊ファイバー $F_0$ が $\mathcal G$ の完全自己同値から誘導されるならば、
$F_A$ は $\mathcal G_A$ の完全自己同値から誘導される。
この自己同値は noetherian 対象および $A$ 上平坦な対象を保つ。
\end{lemma}
\begin{proof}
完全な証明は相対モジュライ論文 \cite[\texttt{artin_foundation}]{QuadricModuli2026} に置く。
\end{proof}

\begin{lemma}[Artin 基底での豊富性の持上げ]
\label{lem:artin-ample-lift}
$F_A$ を前補題の完全自己同値とし、$\mathcal O_A$ を $\mathcal O$ の flat な持上げとする。
$(\mathcal O,F_0)$ が AZ ample であり、全ての $N\in\mathcal C$ について
\[
 \operatorname{Ext}^{j}_{\mathcal C}(F_0^{-n}\mathcal O,N)=0
 \quad(j>0,\ n\gg0)
\]
と仮定する。このとき $(\mathcal O_A,F_A)$ は
$\mathcal C_A=(\mathcal G_A)^{\mathrm{noeth}}$ 上 AZ ample である。
\end{lemma}
\begin{proof}
完全な証明は相対モジュライ論文 \cite[\texttt{artin_foundation}]{QuadricModuli2026} に置く。
\end{proof}

\begin{proposition}[全 Abel 変形の canonical 反 Serre 環]
\label{prop:artin-canonical-antiserre}
任意の flat pointed Abel 変形 $(\mathcal C_A,\mathcal O_A)$ について、
relative compact 圏の反 Serre 関手 $t_A=S_A^{-1}[2]$ は
$\mathcal C_A$ の完全自己同値に延長し、$(\mathcal O_A,t_A)$ は AZ ample である。
さらに
\[
 R_A^{\mathrm{can}}
 =\bigoplus_{n\ge0}\operatorname{Hom}(\mathcal O_A,t_A^n\mathcal O_A)
\]
は全次数自由、階数 $(2n+1)^2$ の二次環である。
生成元の枠と次数線束の標準的な扱いによって $[X_3/\operatorname{PGL}_9]$ の
Artin 変形を与え、
\[
 (\mathcal C_A,\mathcal O_A)
 \simeq(\operatorname{qgr}R_A^{\mathrm{can}},\pi R_A^{\mathrm{can}}).
\]
この構成は Artin 基底変更と両立する。
\end{proposition}
\begin{proof}
完全な証明は相対モジュライ論文 \cite[\texttt{artin_foundation}]{QuadricModuli2026} に置く。
\end{proof}

\begin{proposition}[Artin 関係式族は表示された全軌道環である]
\label{prop:artin-presentation-full-orbit}
命題~\ref{prop:artin-relations-flat-abelian} の $R_A$ について、
$\mathcal O_A=\pi R_A$、$F_A=(1)$ と置くと、
$(\mathcal O_A,F_A)$ は AZ ample であり、自然な次数付き射
\[
 R_A\longrightarrow
 \bigoplus_{n\ge0}\operatorname{Hom}(\mathcal O_A,F_A^n\mathcal O_A)
\]
は全次数で同型である。
\end{proposition}
\begin{proof}
完全な証明は相対モジュライ論文 \cite[\texttt{artin_foundation}]{QuadricModuli2026} に置く。
\end{proof}

\begin{corollary}[形式的な偏極忘却の分裂]
\label{cor:artin-polarization-retraction}
$x\in[X_3/\operatorname{PGL}_9](k)$ を $\HH$ の反 Serre 環の点とする。
Artin 基底上で、関係式から pointed Abel 圏へ移る変形写像
\[
 \mathrm{Def}_{[X_3/\operatorname{PGL}_9],x}
 \longrightarrow\mathrm{Def}_{\mathrm{ab},(\mathcal C,\mathcal O)}
\]
は、canonical 反 Serre 環を取る右逆を持つ。
従って一次変形空間上の写像は分裂全射である。さらに標準的な分裂完全列
\[
 0\longrightarrow\operatorname{HH}^{1}(\Lambda)
 \longrightarrow T_x[X_3/\operatorname{PGL}_9]
 \longrightarrow\operatorname{HH}^{2}(\Lambda)
 \longrightarrow0
 \tag{Artin-pol}\label{eq:artin-polarization-tangent-split}
\]
を得る。ここで中央項は接複体の $H^0$、すなわち閉ファイバーの同定を保つ
一次変形の同型類のベクトル空間である。
二根対一致点ではその核は一次元で、
命題~\ref{prop:antiserre-polarization-tangent} の Zhang 方向が生成する。
\end{corollary}
\begin{proof}
完全な証明は相対モジュライ論文 \cite[\texttt{artin_foundation}]{QuadricModuli2026} に置く。
\end{proof}

\begin{corollary}[Abel 変形の無障害性と三基準点の接次元]
\label{cor:artin-abelian-unobstructed}
$\HH$ の任意の対象の pointed Abel 変形は、Artin 基底の小拡大に対して無障害である。
また
\[
 \dim T_x[X_3/\operatorname{PGL}_9]
 =3+2\dim\operatorname{HH}^{1}(\Lambda).
\]
既存の noncubic 三基準点では、Abel 変形の接次元は $3,3,4$、
偏極を保持した関係式変形の接次元は $3,3,5$ となる。
\end{corollary}
\begin{proof}
完全な証明は相対モジュライ論文 \cite[\texttt{artin_foundation}]{QuadricModuli2026} に置く。
\end{proof}

\begin{corollary}[$\operatorname{HH}^{1}=0$ の点での全 Artin 偏極認識]
\label{cor:artin-rigid-polarization}
$\operatorname{HH}^{1}(\Lambda)=0$ と仮定する。
$x$ を閉ファイバーとする任意の Artin $X_3$-変形について、
表示された shift $F_A=(1)$ は、その pointed Abel 圏の canonical な反 Serre 関手
$t_A=S_A^{-1}[2]$ と自然同型である。
従って偏極忘却は、Artin 変形の同型類の関手上で canonical 反 Serre 環の構成と互いに逆になる。
特に noncubic の二つの根対が異なる全ての点で、この結論が成立する。
根を一つ共有する場合も含む。
\end{corollary}
\begin{proof}
完全な証明は相対モジュライ論文 \cite[\texttt{artin_foundation}]{QuadricModuli2026} に置く。
\end{proof}

\begin{corollary}[非一致根対での関係式空間の滑らかさ]
\label{cor:artin-noncubic-relations-smooth}
$x\in[X_3/\operatorname{PGL}_9](k)$ が、二つの根対の異なる noncubic 対象の
反 Serre 環を表すとする。根を一つ共有する場合も含める。
このとき $[X_3/\operatorname{PGL}_9]$ は $x$ で滑らか、局所次元は三である。
生成元の枠を選んだ対応点 $J\in X_3(k)$ では、$X_3$ は滑らかで局所次元は八十三である。
\end{corollary}
\begin{proof}
完全な証明は相対モジュライ論文 \cite[\texttt{artin_foundation}]{QuadricModuli2026} に置く。
\end{proof}

この滑らかさは $X_3$ とその商についての結論である。
平方根と AS4 被覆の相対存在条件を課す $X_{\mathrm{as}}$ の開性や滑らかさを
結論してはいない。

\begin{remark}
この分裂は Artin 変形の比較である。
canonical な反 Serre 偏極の存在は、適合平方根と AS4 被覆の存在とは異なる。
従って全 Abel 変形が $\mathfrak H_2$ の変形であることや、
$X_{\mathrm{as}}$ の代数性・開性を結論していない。
また、四つの foundation 対象の像だけを比べても偏極は認識できない。
それらは例外対象として一意に持ち上がるため、Zhang 方向でも各像は同型になる。
関手を比較するには、有限 foundation 代数の射への作用との整合も必要である。
\end{remark}

\subsection{Artin 上の平方根・被覆の相対族への昇格}
\label{subsec:artin-relative-promotion}

Artin 基底で対象を作っただけでは、定義~\ref{def:relative-h2} が要求する
全有限型基底変更上の圏の系を確認したことにはならない。
次の命題では、既に平方根と AS 被覆が与えられた場合について、その比較を閉じる。
任意の Abel 圏に対する非平坦基底変更の導来降下定理は用いない。

\begin{proposition}[Artin 上の適合表示から相対 $\mathfrak H_2$ へ]
\label{prop:artin-witness-relative-promotion}
$(A,\mathfrak m,k)$ を有限次元可換局所 Artin $k$-代数とする。
$(\mathcal C_A,\mathcal O_A)$ を $(\mathcal C,\mathcal O)\in\HH$ の
flat pointed Abel 変形とし、次のデータが与えられたと仮定する。
\begin{enumerate}
\item 適合する体上の平方根 $s_0$ の $A$-線形 Morita 持上げ $s_A$。
  有限 foundation の Morita 表示において
  $s_A^2\simeq S_A^{-1}[2]$ という双加群の同型を持つ。
\item 全軌道環の同型
  \[
   B_A:=\bigoplus_{n\ge0}
     \operatorname{Hom}(\mathcal O_A,s_A^n\mathcal O_A)
   \simeq\widetilde B_A/(g_A).
  \]
  ここで $\widetilde B_A$ は各次数有限自由な相対 quadratic AS4 環であり、
  $g_A$ は二次の左右正則正規元、商も全次数 $A$-平坦である。
\end{enumerate}
このとき $s_A$ は heart 上完全で $(\mathcal O_A,s_A)$ は AZ ample である。
さらに、有限型 $A$-代数 $A'$ に対して
\[
 \mathcal C_{A'}=
 \bigl((\operatorname{Ind}\mathcal C_A)_{A'}\bigr)^{\mathrm{noeth}}
 \simeq \operatorname{qgr}(B_A\otimes_A A')
\]
と置き、$\mathcal O_{A'}=A'\otimes_A\mathcal O_A$ とすれば、
これらは定義~\ref{def:relative-h2} の相対族を与える。
従って $(\mathcal C_A,\mathcal O_A)$ は $\mathfrak H_2(A)$ の対象となる。
平方根と被覆はこの相対族の存在を証明するための補助データであり、標識に加えない。
\end{proposition}
\begin{proof}
完全な証明は相対モジュライ論文 \cite[\texttt{artin_foundation}]{QuadricModuli2026} に置く。
\end{proof}

\begin{remark}
本命題は平方根・AS 被覆の\emph{存在}を新たに作っていない。
Artin 基底上でそれらを構成したとき、任意の有限型基底変更を含む相対族の条件を
追加の未証明仮定として残さずに確認するための命題である。
一般の $X_3$-変形に対して AS 被覆が存在するという主張は含まない。
\end{remark}
